% =====================================================================
%  Thème 5 — EXERCICES — Niveau FACILE
% =====================================================================
\documentclass[11pt,a4paper]{article}
\usepackage[utf8]{inputenc}
\usepackage[T1]{fontenc}
\usepackage{lmodern}
\usepackage[french]{babel}
\usepackage{amsmath,amssymb,mathtools}
\usepackage{icomma}
\usepackage{xcolor}
\usepackage[most]{tcolorbox}
\usepackage{enumitem}
\usepackage{array}
\usepackage{booktabs}
\usepackage{fancyhdr}
\usepackage[hidelinks]{hyperref}
\usepackage[a4paper,margin=2.2cm,headheight=26pt,headsep=10pt]{geometry}

\definecolor{primary}{RGB}{223,87,99}
\definecolor{primarydark}{RGB}{150,42,55}

\tcbset{boxbase/.style={enhanced,breakable,boxrule=0.9pt,arc=2.5pt,
   left=3mm,right=3mm,top=2.3mm,bottom=2.3mm,
   fonttitle=\bfseries,coltitle=white,toptitle=0.6mm,bottomtitle=0.6mm}}
\newtcolorbox[auto counter]{exercice}[1][]{boxbase,colback=primary!4,colframe=primary,
  title={\textsc{Exercice}~\thetcbcounter\ifx\relax#1\relax\relax\else\textmd{ --- \itshape #1}\fi}}

\pagestyle{fancy}
\fancyhf{}
\fancyhead[L]{\small\textcolor{primarydark}{\textbf{Fiche 6} — Les limites}}
\fancyhead[R]{\small\textcolor{primarydark}{\textsc{Thème 5 — Facile}}}
\fancyfoot[C]{\small\thepage}
\renewcommand{\headrulewidth}{0.8pt}
\renewcommand{\headrule}{\hbox to\headwidth{\color{primary}\leaders\hrule height \headrulewidth\hfill}}

\newcommand{\R}{\mathbb{R}}

\begin{document}

\begin{center}
{\Large\bfseries\color{primarydark}Thème 5 — L'Hospital \& limites trigonométriques}\\[2pt]
{\large\color{primary}Exercices — Niveau Facile}
\end{center}
\vspace{4pt}
{\small\itshape But : reconnaître une forme $\frac00$ ou $\frac{\infty}{\infty}$, appliquer l'Hospital une
fois, et utiliser $\frac{\sin x}{x}\to 1$, $\frac{\tan x}{x}\to 1$.}
\vspace{6pt}

\begin{exercice}[l'Hospital est-il applicable ?]
Pour chaque limite, dire si la forme directe est $\frac00$, $\frac{\infty}{\infty}$, ou une forme
\emph{déterminée} (auquel cas l'Hospital ne s'applique pas).
\begin{enumerate}[label=\alph*),leftmargin=2em,itemsep=3pt]
  \item $\displaystyle\lim_{x\to 0}\frac{\sin x}{x}$ \qquad b) $\displaystyle\lim_{x\to 0}\frac{\cos x}{x+1}$
  \item $\displaystyle\lim_{x\to+\infty}\frac{x}{e^x}$ \qquad d) $\displaystyle\lim_{x\to 2}\frac{x^2}{x-1}$
\end{enumerate}
\end{exercice}

\begin{exercice}[l'Hospital, une fois]
Chaque forme est $\frac00$. Appliquer l'Hospital (dériver haut et bas séparément).
\begin{enumerate}[label=\alph*),leftmargin=2em,itemsep=3pt]
  \item $\displaystyle\lim_{x\to 0}\frac{e^x-1}{x}$
  \item $\displaystyle\lim_{x\to 0}\frac{\sin x}{x}$
  \item $\displaystyle\lim_{x\to 1}\frac{\ln x}{x-1}$
\end{enumerate}
\end{exercice}

\begin{exercice}[limites trigonométriques directes]
Donner la valeur (résultats de référence).
\begin{enumerate}[label=\alph*),leftmargin=2em,itemsep=3pt]
  \item $\displaystyle\lim_{x\to 0}\frac{\sin x}{x}$ \qquad b) $\displaystyle\lim_{x\to 0}\frac{\tan x}{x}$
  \item $\displaystyle\lim_{x\to 0}\frac{\sin 2x}{2x}$
\end{enumerate}
\end{exercice}

\begin{exercice}[reconstituer $\frac{\sin(\square)}{\square}$]
Calculer en faisant apparaître une forme de référence.
\begin{enumerate}[label=\alph*),leftmargin=2em,itemsep=3pt]
  \item $\displaystyle\lim_{x\to 0}\frac{\sin 2x}{3x}$
  \item $\displaystyle\lim_{x\to 0}\frac{\sin 3x}{x}$
\end{enumerate}
\end{exercice}

\begin{exercice}[l'Hospital sur un quotient polynomial]
La forme est $\frac00$. Calculer par l'Hospital, puis vérifier par factorisation.
\[ \lim_{x\to 1}\frac{x^2-1}{x-1}. \]
\end{exercice}

\end{document}
